
\subsection{Deployment}
\label{sec:deployment}

\subsubsection{Embodiment-specific fine-tuning}
\label{sec:finetuning}



\paragraph{Shared recipe.}
Embodiment-specific fine-tuning starts from the post-trained \molmoacttpos checkpoint and keeps the same VLM--action-expert architecture described in Sec.~\ref{sec:posttraining}. We continue to co-train the discrete autoregressive action output and the continuous flow-matching action expert, use the same 32-dimensional padded action width, and keep the same learning rates as post-training: \(5\times10^{-6}\) for the vision encoder and connector, \(1\times10^{-5}\) for the language model, and \(5\times10^{-5}\) for the action expert. We also keep single resized crops, packed sequences, setup/control descriptors, discrete state tokens, and the same action-token vocabulary, using the shared prompt, image-augmentation, and packing implementation described in \autoref{supp:training}.

The fine-tuning stage differs from post-training in four ways. First, each run is robot-only: we do not include the multimodal VLM mixture or a separate VLM dataloader. Second, we increase the number of sampled flow times from 4 to 8 per action chunk, which gives the action expert denser supervision along the same flow trajectory. Third, we do not use knowledge insulation during fine-tuning: gradients from the flow loss are allowed to update the VLM through the action-expert conditioning path, since we did not observe a consistent performance gain from detaching this path at this stage. Finally, for the main embodiment checkpoints below, we do not tune the added-token input embeddings. As in standard VLM fine-tuning, we assume that the token embeddings learned during pre-training and post-training are already well placed, and tune the language-model output head and final normalization instead.

\paragraph{Bimanual YAM.}
The \molmoacttyam checkpoint is fine-tuned on the bimanual YAM mixture. The camera order is fixed as top, left, and right, matching the data collection setup. We use annotated language instructions, absolute joint-pose actions, and a 30-step action chunk, corresponding to the 30 Hz control rate of the dataset. We train with a sequence length of 2100, global batch size 128, and 100K updates on 64 H100 GPUs, for roughly 2,304 GPU hours.

\paragraph{DROID.}
The \molmoacttdroid checkpoint is fine-tuned on the filtered DROID mixture. DROID provides two exterior cameras and one wrist camera. For DROID-style two-view evaluation, we use a fixed exterior-then-wrist ordering; when training variants expose both exterior choices as alternatives, the loader samples one exterior camera and pairs it with the wrist camera during training. We use absolute joint-pose actions with a 15-step action chunk, matching DROID's 15 Hz control rate. We do not use the additional language annotations in this fine-tune, to keep the comparison with prior DROID-trained models fair. We train with a sequence length of 2100, global batch size 64, and 100K updates on 32 H100 GPUs, for roughly 1,152 GPU hours.

\paragraph{SO-100/101.}
The \molmoacttso checkpoint is fine-tuned on the SO-100/101 mixture. Because this data is aggregated from internet demonstrations with diverse and inconsistent camera layouts, we randomize the input camera order at the episode level rather than imposing a fixed view naming convention. We use annotated language instructions, absolute joint-pose actions, and a 30-step action chunk for the 30 Hz control rate. We train with a sequence length of 2100, global batch size 64, and 100K updates on 32 H100 GPUs, for roughly 1,152 GPU hours.

\paragraph{LIBERO.}
The \molmoacttlib checkpoint is fine-tuned on the full LIBERO training mixture, combining the Spatial, Object, Goal, and Long suites rather than training a separate model for each suite. The camera order is fixed as front view followed by wrist view. We do not use annotated language instructions. LIBERO uses relative end-effector control at 10 Hz, so we train with a 10-step action chunk. We use a sequence length of 2100, global batch size 64, and train for 50K updates on 32 H100 GPUs, for roughly 1,152 GPU hours; for evaluation, we select the best-performing checkpoint, which occurs at 40K updates.

\paragraph{Other evaluation fine-tunes.}
For smaller task- or benchmark-specific fine-tunes used in downstream evaluation, we follow the same recipe unless otherwise noted. These runs use fixed metadata camera order rather than camera-order randomization, no language annotations, 2100-token sequences, packing, 8 sampled flow times, and robot-only data. The action horizon is set to the dataset control frequency, and the control representation follows the dataset, either joint pose or end-effector pose. Real-world evaluation runs use 8 H100 GPUs, global batch size 16, and 50K updates, and we choose the best-performing checkpoint for evaluation.

\subsubsection{Inference optimization}
\label{sec:inference}
For the \molmoactt model, inference is dominated by the continuous action expert, which integrates a fixed-step flow-matching trajectory conditioned on the VLM context. Within one action chunk, this context is invariant across flow steps, while only the noisy action state and flow time evolve. We therefore cache reusable action-expert intermediates, including context-dependent cross-attention states and fixed position-dependent terms, and reuse them throughout the flow loop. We further capture the fixed-shape flow loop with CUDA Graphs, reducing Python and kernel-launch overhead during repeated action generation.
